\chapter{Hito 5: Construcción de la Capa Staging en DuckDB}

\section{Objetivo}

El objetivo de este hito es construir la primera capa física del Data Warehouse utilizando DuckDB.

La capa \texttt{staging} será el punto de entrada de todos los datos al sistema. En esta capa almacenaremos los datos provenientes de archivos CSV en una forma cercana a la fuente original.

Al finalizar este hito el estudiante será capaz de:

\begin{itemize}
\item Crear schemas en DuckDB.
\item Crear tablas de aterrizaje para datos crudos.
\item Ejecutar scripts SQL desde Python.
\item Preparar la base local para recibir datos desde archivos CSV.
\item Mantener la arquitectura \texttt{staging -> core -> mart}.
\end{itemize}

\section{Resultado esperado}

Al finalizar este hito deberán existir tres schemas:

\begin{verbatim}
staging
core
mart
\end{verbatim}

y dos tablas en la capa \texttt{staging}:

\begin{verbatim}
staging.tipo_cambio_raw

staging.balance_bancario_raw
\end{verbatim}

\section{Paso 1: Crear script de schemas}

Crear el archivo:

\begin{verbatim}
sql/01_create_schemas.sql
\end{verbatim}

Contenido:

\begin{lstlisting}[language=SQL]
CREATE SCHEMA IF NOT EXISTS staging;

CREATE SCHEMA IF NOT EXISTS core;

CREATE SCHEMA IF NOT EXISTS mart;
\end{lstlisting}

\section{Paso 2: Crear script para ejecutar SQL}

Crear o verificar el archivo:

\begin{verbatim}
etl/run_sql_file.py
\end{verbatim}

Contenido:

\begin{lstlisting}[language=Python]
from db import get_connection

def run_sql_file(path):

```
conn = get_connection()

with open(path, "r", encoding="utf-8") as file:
    sql = file.read()

conn.execute(sql)

conn.close()
```

if **name** == "**main**":

```
run_sql_file(
    "sql/01_create_schemas.sql"
)

print("Script SQL ejecutado correctamente.")
```

\end{lstlisting}

\section{Paso 3: Ejecutar creación de schemas}

Desde la terminal:

\begin{lstlisting}[language=bash]
python etl/run_sql_file.py
\end{lstlisting}

Salida esperada:

\begin{verbatim}
Script SQL ejecutado correctamente.
\end{verbatim}

\section{Paso 4: Verificar schemas}

Crear el archivo:

\begin{verbatim}
sql/check_schemas.sql
\end{verbatim}

Contenido:

\begin{lstlisting}[language=SQL]
SELECT schema_name
FROM information_schema.schemata
WHERE schema_name IN
(
'staging',
'core',
'mart'
)
ORDER BY schema_name;
\end{lstlisting}

Modificar temporalmente \texttt{etl/run_sql_file.py} para ejecutar:

\begin{lstlisting}[language=Python]
run_sql_file(
"sql/check_schemas.sql"
)
\end{lstlisting}

Resultado esperado:

\begin{verbatim}
core
mart
staging
\end{verbatim}

\section{Paso 5: Crear script de tablas staging}

Crear:

\begin{verbatim}
sql/02_staging_ddl.sql
\end{verbatim}

\section{Paso 6: Crear tabla tipo_cambio_raw}

Agregar:

\begin{lstlisting}[language=SQL]
CREATE TABLE IF NOT EXISTS staging.tipo_cambio_raw
(
fecha TEXT,

```
moneda TEXT,

tasa_compra TEXT,

tasa_venta TEXT,

archivo_origen TEXT,

cargado_en TIMESTAMP
```

);
\end{lstlisting}

\section{Paso 7: Crear tabla balance_bancario_raw}

Agregar:

\begin{lstlisting}[language=SQL]
CREATE TABLE IF NOT EXISTS staging.balance_bancario_raw
(
periodo TEXT,

```
entidad_cod TEXT,

entidad_nombre TEXT,

tipo_entidad TEXT,

activos_totales TEXT,

cartera_creditos TEXT,

depositos TEXT,

patrimonio TEXT,

resultado_neto TEXT,

depositos_me TEXT,

cartera_me TEXT,

archivo_origen TEXT,

cargado_en TIMESTAMP
```

);
\end{lstlisting}

\section{Nota sobre tipos de datos}

En \texttt{staging} usaremos columnas tipo \texttt{TEXT}.

La idea es cargar primero los datos crudos y validar después.

Esta decisión permite:

\begin{itemize}
\item Evitar fallos tempranos por formatos incorrectos.
\item Conservar el dato original.
\item Separar carga de transformación.
\item Diagnosticar errores con mayor facilidad.
\end{itemize}

\section{Paso 8: Ejecutar DDL de staging}

Modificar temporalmente:

\begin{verbatim}
etl/run_sql_file.py
\end{verbatim}

para ejecutar:

\begin{lstlisting}[language=Python]
run_sql_file(
"sql/02_staging_ddl.sql"
)
\end{lstlisting}

Luego ejecutar:

\begin{lstlisting}[language=bash]
python etl/run_sql_file.py
\end{lstlisting}

\section{Paso 9: Verificar tablas creadas}

Crear:

\begin{verbatim}
sql/check_staging_tables.sql
\end{verbatim}

Contenido:

\begin{lstlisting}[language=SQL]
SELECT table_schema,
table_name
FROM information_schema.tables
WHERE table_schema = 'staging'
ORDER BY table_name;
\end{lstlisting}

Resultado esperado:

\begin{verbatim}
staging    balance_bancario_raw
staging    tipo_cambio_raw
\end{verbatim}

\section{Paso 10: Crear script general para inicializar la base}

Crear:

\begin{verbatim}
etl/init_database.py
\end{verbatim}

Contenido:

\begin{lstlisting}[language=Python]
from db import get_connection

def run_sql(path, conn):

```
with open(path, "r", encoding="utf-8") as file:
    sql = file.read()

conn.execute(sql)
```

def main():

```
conn = get_connection()

run_sql(
    "sql/01_create_schemas.sql",
    conn
)

run_sql(
    "sql/02_staging_ddl.sql",
    conn
)

conn.close()

print("Base DuckDB inicializada correctamente.")
```

if **name** == "**main**":
main()
\end{lstlisting}

\section{Paso 11: Ejecutar inicialización completa}

Desde la terminal:

\begin{lstlisting}[language=bash]
python etl/init_database.py
\end{lstlisting}

Salida esperada:

\begin{verbatim}
Base DuckDB inicializada correctamente.
\end{verbatim}

\section{Paso 12: Crear consulta de monitoreo}

Crear:

\begin{verbatim}
sql/monitor_staging.sql
\end{verbatim}

Contenido:

\begin{lstlisting}[language=SQL]
SELECT
'tipo_cambio_raw' AS tabla,
COUNT(*) AS total_filas
FROM staging.tipo_cambio_raw

UNION ALL

SELECT
'balance_bancario_raw' AS tabla,
COUNT(*) AS total_filas
FROM staging.balance_bancario_raw;
\end{lstlisting}

Inicialmente ambas tablas deberán tener cero filas.

\section{Buenas prácticas}

Durante este hito seguiremos estas reglas:

\begin{enumerate}
\item La base DuckDB vive en \texttt{data/warehouse}.
\item Los schemas separan las capas del Data Warehouse.
\item La capa \texttt{staging} recibe datos crudos.
\item Las transformaciones ocurrirán en \texttt{core}.
\item Las vistas para análisis se construirán en \texttt{mart}.
\end{enumerate}

\section{Checklist de validación}

Antes de continuar verificar:

\begin{itemize}
\item[$\Box$] Archivo \texttt{macrofinanzas.duckdb} existe.
\item[$\Box$] Schema \texttt{staging} creado.
\item[$\Box$] Schema \texttt{core} creado.
\item[$\Box$] Schema \texttt{mart} creado.
\item[$\Box$] Tabla \texttt{tipo_cambio_raw} creada.
\item[$\Box$] Tabla \texttt{balance_bancario_raw} creada.
\item[$\Box$] Script \texttt{init_database.py} ejecutado.
\item[$\Box$] Consulta de monitoreo creada.
\end{itemize}

\section{Entregable}

Al finalizar este hito el estudiante deberá disponer de:

\begin{itemize}
\item Base DuckDB inicializada.
\item Schemas creados.
\item Tablas staging creadas.
\item Script de inicialización.
\item Capa Bronze lista para recibir datos.
\end{itemize}

\section{Próximo hito}

En el Hito 6 construiremos el proceso ETL de carga.

Aprenderemos a:

\begin{itemize}
\item Leer archivos CSV.
\item Cargar datos hacia DuckDB.
\item Limpiar la capa staging antes de cada carga.
\item Automatizar el flujo de ingesta.
\item Verificar conteos de registros cargados.
\end{itemize}
